\partpage{B}{Two Marks Questions}
\chap{1}{Electric Charges and Fields}

\Q{1}{Define electric field at a point. What is meant by `source charge'?}{MQP-1, Q.21}
\RA{Electric field at a point is the electrostatic force experienced by a unit positive (test) charge placed at that point.}{1}
\R{The charge which produces the electric field is called the source charge.}{1}

\Q{2}{Sketch the electric field lines for (a) a positive point charge and (b) an electric dipole.}{MQP-2, Q.21}
\nfig{1_14}{0.80}{1.14 --- Field lines due to (a) a positive point charge and (d) an electric dipole (Physics Part I, Ch.\ 1)}
\R{(Arrows on the field lines are essential: outward for a positive charge; from $+q$ to $-q$ for a dipole.)}{1+1}

\Q{3}{Define electric flux through an area element. Mention its SI unit.}{MQP-3, Q.21}
\RA{Electric flux through an area element is the product of the electric field and the component of the area element perpendicular to the field (the number of field lines crossing the area).}{1}
\R{SI unit: N\,C$^{-1}$\,m$^2$ \ OR \ V\,m \ OR \ J\,C$^{-1}$\,m \ OR \ N\,A$^{-1}$\,s$^{-1}$\,m$^2$}{1}
\src{SI unit as in KSEAB Scheme of Evaluation, II PUC Annual Exam March 2020, Q.1}

\Q{4}{Define the term `linear charge density'. Mention its SI unit.}{MQP-4, Q.21}
\RA{Linear charge density is the charge per unit length of a line charge.}{1}
\R{SI unit: coulomb per metre (C\,m$^{-1}$).}{1}

\Q{5}{Mention two factors on which electric field at a point due to a point charge depends.}{MQP-5, Q.21}
\RA{(i) Magnitude of the source charge $q$ ($E\propto q$).}{1}
\R{(ii) Distance $r$ of the point from the charge ($E\propto 1/r^2$).}{1}
\R{(iii) Permittivity of the medium. \hfill(Any two)}{}

\chap{2}{Electrostatic Potential and Capacitance}

\Q{6}{The electrostatic potential energy of a system of two like charges is positive. Explain.}{MQP-1, Q.22}
\RA{Potential energy of two point charges: $U=\dfrac{1}{4\pi\varepsilon_0}\dfrac{q_1q_2}{r}$.}{1}
\R{For like charges $q_1q_2>0$, so $U>0$. Positive work has to be done against the force of repulsion to bring them closer from infinity; this work is stored as positive potential energy.}{1}

\Q{7}{What are polar molecules? Give an example.}{MQP-2, Q.22}
\RA{The molecules in which the centres of positive charges and negative charges do not coincide are called polar molecules. They have a permanent (intrinsic) dipole moment.}{1}
\R{Example: water (H$_2$O), HCl.}{1}
\src{KSEAB Scheme of Evaluation, II PUC Supplementary Exam August 2021, Q.12 (definition)}

\Q{8}{What is meant by equipotential surface? Draw equipotential surfaces for an electric dipole.}{MQP-3, Q.22}
\RA{An equipotential surface is a surface with a constant value of potential at all points on the surface.}{1}
\src{KSEAB Scheme of Evaluation, II PUC Annual Exam March 2018, Q.1}
\nfig{2_11a}{0.40}{2.11 (a) --- Equipotential surfaces for a dipole (Physics Part I, Ch.\ 2)}
\R{}{1}

\Q{9}{Write two properties of equipotential surfaces.}{MQP-4, Q.22}
\RA{(i) For any charge configuration, the equipotential surface through a point is normal to the electric field at that point.}{}
\R{(ii) Work done to move a charge on an equipotential surface is zero.}{2}
\R{(iii) The equipotential surfaces corresponding to different potentials are (a) very close in case of a strong field and (b) far apart in case of a weak field.}{}
\R{(iv) Two equipotential surfaces do not intersect. \hfill(Any two)}{}
\src{KSEAB Scheme of Evaluation, II PUC Supplementary Exam May 2023, Q.21}

\Q{10}{The electrostatic potential energy of a system of two like charges decreases with increase in the distance of separation. Explain the statement using suitable expression.}{MQP-5, Q.22}
\RA{$U=\dfrac{1}{4\pi\varepsilon_0}\dfrac{q_1q_2}{r}$}{1}
\R{For like charges $q_1q_2$ is positive, so $U\propto\dfrac1r$. As the separation $r$ increases, $U$ decreases (the repulsive force does work as the charges move apart).}{1}

\chap{3}{Current Electricity}

\Q{11}{Name the mobile charge carriers in metals. Write the expression for mobility in terms of average collision time (relaxation time) for charge carriers.}{MQP-1, Q.23}
\RA{Mobile charge carriers in metals: free electrons.}{1}
\R{Mobility: $\mu=\dfrac{|v_d|}{E}=\dfrac{e\tau}{m}$, where $\tau$ = relaxation time, $m$ = mass and $e$ = charge of the electron.}{1}

\Q{12}{State and explain Ohm's law.}{MQP-2, Q.23}
\RA{\textbf{Ohm's law:} The electric current ($I$) flowing through a substance (conductor) is directly proportional to the voltage (potential difference, $V$) across its ends, provided the temperature and other physical conditions remain constant (same).}{1}
\R{If $I$ is the current through the conductor and $V$ the potential difference between its ends, $V\propto I$, i.e.\ $V=IR$, where the constant $R$ is the resistance of the conductor.}{1}
\src{KSEAB Scheme of Evaluation, II PUC Supplementary Exam August 2021, Q.21 and Annual Exam March 2017, Q.13}

\Q{13}{Write two limitations of Ohm's law.}{MQP-3, Q.23}
\RA{(1) $V$ ceases to be proportional to $I$.}{1}
\R{(2) Semiconductor diodes do not obey Ohm's law.}{+}
\R{(3) For some materials like GaAs, the relation between $V$ and $I$ is not unique.}{1}
\R{(4) Ohm's law is not applicable at very high and very low temperatures.\\ (Any two) (Any other correct limitation should be considered)}{}
\src{KSEAB Scheme of Evaluation, II PUC Exam-1 March 2025, Q.22}

\Q{14}{Draw the graphs to show the variation of resistivity with temperature in case of (a) conductors and (b) semiconductors.}{MQP-4, Q.23}
\ofig{og13}{0.62}{(a) Conductor (copper) \qquad\qquad (b) Semiconductor (silicon)}
\R{}{1+1}
\src{KSEAB Scheme of Evaluation, II PUC Exam-2 2026, Q.23 (graphs reproduced from the scheme)}
\note{(For conductors resistivity increases with temperature; for semiconductors it decreases.)}

\Q{15}{State and explain Kirchhoff's loop rule.}{MQP-5, Q.23}
\RA{\textbf{Loop rule:} In an electrical network, the algebraic sum of changes in potential around any closed loop involving resistors and cells in the loop is zero.}{1}
\src{KSEAB Scheme of Evaluation, II PUC Exam-2 April/May 2025, Q.23}
\R{\textbf{Explanation:} Electric potential depends only on position, so on returning to the starting point the total change in potential is zero ($\sum\Delta V=0$). The rule is based on the law of conservation of energy. E.g.\ for a loop with a cell $\varepsilon$ and resistors $R_1$, $R_2$: $\varepsilon-IR_1-IR_2=0$.}{1}

\chap{4}{Moving Charges and Magnetism}

\Q{16}{What is the nature of the force between two parallel conductors carrying currents in the (a) same direction and (b) opposite directions?}{MQP-1, Q.24}
\RA{(a) Same direction: \textbf{attractive}.}{1}
\R{(b) Opposite directions: \textbf{repulsive}.}{1}

\Q{17}{A current loop of magnetic dipole moment 10 A\,m$^2$ is suspended freely in a uniform magnetic field of magnitude $5\times10^{-3}$\,T. The plane of the coil is in the direction of the magnetic field. Calculate the torque on the current loop.}{MQP-2, Q.24}
\RA{Torque: $\tau=mB\sin\theta$. Since the plane of the coil is along $\vb B$, $\vb m\perp\vb B$, so $\theta=90^\circ$.}{1}
\R{$\tau=10\times5\times10^{-3}\times\sin90^\circ \Rightarrow \boxed{\tau=5\times10^{-2}\ \text{N\,m}}$}{1}

\Q{18}{Give the vector form of Biot--Savart law and explain the terms.}{MQP-3, Q.24}
\RA{$d\vb B=\dfrac{\mu_0}{4\pi}\,\dfrac{I\,d\vb l\times\vb r}{r^3}$ \quad OR \quad $d\vb B=\dfrac{\mu_0}{4\pi}\,\dfrac{I\,d\vb l\times\hat r}{r^2}$}{1}
\R{$d\vb B$ = magnetic field at the point due to the current element $I\,d\vb l$; $\vb r$ = position vector of the point from the element ($\hat r$ its unit vector); $\mu_0$ = permeability of free space ($4\pi\times10^{-7}$\,T\,m\,A$^{-1}$).}{1}
\ofig{og17}{0.30}{Current element $I\,d\vb l$ and the field $d\vb B$ at P (from the KSEAB scheme)}
\src{KSEAB Scheme of Evaluation, II PUC Exam-2 2026, Q.31 (expression and diagram)}

\Q{19}{A straight wire of length 0.5\,m carrying a current of 2\,A is placed in a uniform magnetic field of 0.3\,T. Calculate the maximum magnetic force on the wire.}{MQP-4, Q.24}
\RA{$F=BIl\sin\theta$; it is maximum when $\theta=90^\circ$: $F_{\max}=BIl$}{1}
\R{$F_{\max}=0.3\times2\times0.5 \Rightarrow \boxed{F_{\max}=0.3\ \text{N}}$}{1}

\Q{20}{When is the force on a charged particle moving in a uniform magnetic field (a) maximum and (b) zero?}{MQP-5, Q.24}
\note{The magnitude of the force is $F=qvB\sin\theta$, where $\theta$ is the angle between $\vb v$ and $\vb B$ (KSEAB Supplementary Exam May 2023, Q.22).}
\RA{(a) \textbf{Maximum} ($F=qvB$) when the particle moves perpendicular to the field ($\theta=90^\circ$).}{1}
\R{(b) \textbf{Zero} when the particle moves parallel or antiparallel to the field ($\theta=0^\circ$ or $180^\circ$).}{1}

\chap{6}{Electromagnetic Induction}

\Q{21}{Calculate the magnetic energy stored in a coil of self-inductance 0.05\,H carrying a current of 0.2\,A.}{MQP-1, Q.25}
\RA{$U=\tfrac12LI^2$}{1}
\R{$U=\tfrac12\times0.05\times(0.2)^2 \Rightarrow \boxed{U=1\times10^{-3}\ \text{J}}$}{1}

\Q{22}{Mention two factors on which self-inductance of a solenoid depends.}{MQP-2, Q.25}
\RA{(i) Number of turns per unit length $n$ \quad (ii) Area of cross-section $A$}{1}
\R{(iii) Length $l$ of the solenoid \quad (iv) Relative permeability $\mu_r$ of the core material \hfill(Any two)}{1}

\Q{23}{What is AC generator? What is its working principle?}{MQP-3, Q.25}
\RA{An AC generator is a device which converts mechanical energy into electrical energy (alternating emf).}{1}
\R{\textbf{Principle:} Electromagnetic induction. When a coil is rotated in a magnetic field, the magnetic flux linked with it changes and an emf is induced.}{1}

\Q{24}{State and explain Faraday's law of electromagnetic induction.}{MQP-4, Q.25}
\RA{The magnitude of the induced emf in a circuit is equal to the time rate of change of magnetic flux through the circuit. \textbf{OR} The magnitude of the induced emf in a circuit is directly proportional to the rate of change of magnetic flux linked with the circuit.}{1}
\R{If $d\phi_B$ is the change in magnetic flux in time $dt$, then the induced emf is $\varepsilon=-\dfrac{d\phi_B}{dt}$.}{1}
\src{KSEAB Scheme of Evaluation, II PUC Annual Exam March 2017, Q.15}

\Q{25}{Give the expression for mutual inductance between two co-axial solenoids and explain the terms.}{MQP-5, Q.25}
\RA{$M=\mu_0n_1n_2\pi r_1^2l$}{1}
\R{$n_1$, $n_2$ = number of turns per unit length of the inner and outer solenoids, $r_1$ = radius of the inner solenoid, $l$ = length of the solenoids, $\mu_0$ = permeability of free space.}{1}

\chap{8}{Electromagnetic Waves}

\Q{26}{Name the electromagnetic waves used in: (a) remote switches of TV sets, (b) the laser assisted eye surgery (LASIK).}{MQP-1, Q.26}
\RA{(a) Remote switches of TV sets $\rightarrow$ \textbf{Infra-red rays}}{1}
\R{(b) LASIK surgery $\rightarrow$ \textbf{Ultraviolet rays}}{1}
\src{KSEAB Scheme of Evaluation, II PUC Supplementary Exam September 2020, Q.17 and Exam-1 March 2026, Q.26}

\Q{27}{What is displacement current? Write the mathematical form of Ampere--Maxwell law.}{MQP-2, Q.26}
\RA{The current due to changing electric flux (or electric field) with time is called displacement current.}{1}
\R{Ampere--Maxwell law: $\displaystyle\oint\vb B\cdot d\vb l=\mu_0 i_c+\mu_0\varepsilon_0\frac{d\phi_E}{dt}$}{1}
\src{KSEAB Scheme of Evaluation, II PUC Exam-2 2026, Q.26}

\Q{28}{List two uses of microwaves.}{MQP-3, Q.26}
\RA{(1) They are used in radar systems used in aircraft navigation.}{1}
\R{(2) They are used in speed guns.}{+}
\R{(3) They are used in microwave ovens.\\ (Any two) (Any other relevant use should be considered)}{1}
\src{KSEAB Scheme of Evaluation, II PUC Exam-1 March 2025, Q.25}

\Q{29}{Name the electromagnetic waves generated by: (a) the accelerated motion of charges in conducting wires, (b) the bombardment of a metal target by high energy electrons.}{MQP-4, Q.26}
\RA{(a) Radio waves}{1}
\R{(b) X-rays}{1}

\Q{30}{Mention two applications of X-rays.}{MQP-5, Q.26}
\RA{(a) X-rays are used as a diagnostic tool in medicine.}{1}
\R{(b) To treat certain forms of cancer.}{1}
\src{KSEAB Scheme of Evaluation, II PUC Annual Exam March 2018, Q.15}

\chap{12}{Atoms}

\Q{31}{Give two limitations of Bohr's atomic model.}{MQP-1, Q.27}
\RA{(i) Bohr's model is applicable only to hydrogenic (single electron) atoms. \textbf{OR} Bohr's model is not applicable to atoms with two or more electrons.}{1}
\R{(ii) This model is unable to explain the relative intensities of the frequencies emitted by hydrogenic (hydrogen-like) atoms. \hfill(Any other correct answer)}{1}
\src{KSEAB Scheme of Evaluation, II PUC Exam-2 2026, Q.27}

\Q{32}{Define impact parameter. What is the value of angle of scattering for minimum impact parameter?}{MQP-2, Q.27}
\RA{The impact parameter is the perpendicular distance of the initial velocity vector of the $\alpha$-particle from the centre of the target nucleus.}{1}
\R{Angle of scattering is $180^\circ$ (for zero, i.e.\ minimum, impact parameter the $\alpha$-particle rebounds).}{1}
\src{KSEAB Scheme of Evaluation, II PUC Exam-1 March 2026, Q.27}

\Q{33}{The radius of the innermost electron orbit of a hydrogen atom is $5.3\times10^{-11}$\,m. Determine the radius of the $n=2$ orbit.}{MQP-3, Q.27}
\RA{$r_n=n^2r_1$}{1}
\R{$r_2=2^2\times5.3\times10^{-11}\Rightarrow\boxed{r_2=2.12\times10^{-10}\ \text{m}}$}{1}

\Q{34}{`Stable structure of an atom as given by Rutherford model is not possible according to classical electromagnetic theory'. Explain this statement.}{MQP-4, Q.27}
\RA{According to classical electromagnetic theory, an accelerated charge radiates energy. An electron revolving in a circular orbit is accelerated (centripetal acceleration), so it continuously loses energy.}{1}
\R{Its orbit would shrink and the electron would spiral into the nucleus. Hence the atom cannot be stable (and it would emit a continuous spectrum, contrary to the observed line spectra).}{1}

\Q{35}{Calculate the angular momentum of an electron revolving in the 1st orbit of hydrogen atom.}{MQP-5, Q.27}
\RA{$L=\dfrac{nh}{2\pi}$; for $n=1$, $L=\dfrac{h}{2\pi}$}{1}
\R{$L=\dfrac{6.63\times10^{-34}}{2\times3.14}\Rightarrow\boxed{L=1.055\times10^{-34}\ \text{J\,s}}$}{1}

\chap{14}{Semiconductor Electronics}

\Q{36}{List out two differences between p-type and n-type semiconductors.}{MQP-1, Q.28}
\Ans\par
\begin{center}\small\begin{tabular}{|p{0.38\linewidth}|p{0.38\linewidth}|}\hline
\textbf{p-type semiconductors} & \textbf{n-type semiconductors}\\\hline
1. They are obtained by doping pure semiconductors with trivalent impurity. & 1. They are obtained by doping pure semiconductors with pentavalent impurity.\\\hline
2. Holes are the majority charge carriers. & 2. Electrons are the majority charge carriers.\\\hline
3. Electrons are the minority charge carriers. & 3. Holes are the minority charge carriers.\\\hline
\end{tabular}\hfill 1+1\end{center}
\note{(Any two; any other correct difference must be considered)}
\src{KSEAB Scheme of Evaluation, II PUC Exam-1 March 2026, Q.28}

\Q{37}{Give any two differences between intrinsic and extrinsic semiconductors.}{MQP-2, Q.28}
\Ans\par
\begin{center}\small\begin{tabular}{|p{0.38\linewidth}|p{0.38\linewidth}|}\hline
\textbf{Intrinsic semiconductors} & \textbf{Extrinsic semiconductors}\\\hline
Semiconductor without doping (pure form). & Semiconductor doped with impurity atoms (impure form).\\\hline
Number of free electrons ($n_e$) is equal to the number of holes ($n_h$), $n_e=n_h$. & Number of free electrons and holes are unequal, $n_e\neq n_h$.\\\hline
Conductivity is low. & Conductivity is high.\\\hline
Conductivity depends only on temperature. & Conductivity depends both on temperature and impurity added.\\\hline
Conductivity is equally due to electrons and holes. & Conductivity is mainly due to majority charge carriers.\\\hline
\end{tabular}\hfill 1+1\end{center}
\note{(Any two --- any other correct difference)}
\src{KSEAB Scheme of Evaluation, II PUC Exam-2 2026, Q.28}

\Q{38}{Mention two differences between half wave rectifier and full wave rectifier.}{MQP-3, Q.28}
\Ans\par
\begin{center}\small\begin{tabular}{|p{0.38\linewidth}|p{0.38\linewidth}|}\hline
\textbf{Half-wave rectifier} & \textbf{Full-wave rectifier}\\\hline
Rectifies only one half cycle of the AC input. & Rectifies both half cycles of the AC input.\\\hline
Uses one diode. & Uses two diodes (with a centre-tap transformer).\\\hline
Efficiency is low (about 40.6\%). & Efficiency is high (about 81.2\%).\\\hline
Output has more ripple. & Output has less ripple.\\\hline
\end{tabular}\hfill 1+1\end{center}
\note{(Any two)}

\Q{39}{What are dopants? Name the type of dopants used to obtain p-type semiconductors.}{MQP-4, Q.28}
\RA{Dopants are the impurity atoms deliberately added to a pure semiconductor to increase its conductivity.}{1}
\R{p-type: \textbf{trivalent} (acceptor) impurities, e.g.\ boron, aluminium, indium, gallium.}{1}

\Q{40}{What is meant by energy gap? Give its value for a conductor.}{MQP-5, Q.28}
\RA{The energy difference between the bottom of the conduction band and the top of the valence band is called the energy gap (forbidden gap). No electron energy states exist in this gap.}{1}
\R{For a conductor, the energy gap is \textbf{zero} (the valence and conduction bands overlap).}{1}
\src{Value for a conductor: KSEAB Scheme of Evaluation, II PUC Supplementary Exam May 2023, Q.20}
